\documentclass[a4paper,twoside]{report}
\usepackage[print]{autiwa}

\setcounter{secnumdepth}{2}

\fancyhead[LE]{}

\makeatletter

% macros qui servent à la mise en page, elles ne doivent pas être utilisées directement
\newcommand{\preparationTopSep}{\vspace{.2cm} \hrule height0.25pt width\hsize \vspace{1em}}
\newcommand{\statistiqueTopSep}{\vspace{.3cm} \hrule height1pt width\hsize \nobreak \vskip\parskip \vspace{.3cm}}
\newcommand{\ingredientsTopSep}{\vspace{.4cm} \hrule height0.75pt width\hsize\vspace*{1\p@}\hrule height0.25pt width\hsize \vspace{1em}}
\newcommand{\partStyle}{\bfseries \large}

% redéfinition d'une nouvelle commande de section qui n'a pas de numérotation, est centrée, avec mise en page modifiée. C'est pour afficher le titre des recettes et qu'elles apparaissent dans le menu. J'avais fait avec addcontentline mais la référence n'était pas bonne et pointait sur la page précédente parfois. Bref, pas propre.
\newcommand\nomRecette{\@startsection {section}{6}{\z@}{0ex}{2.3ex }%
	{\reset@font\Huge\bfseries\centering}}

\usepackage{bbding}%to add \FiveStar and \FiveStarOpen
\newcommand{\note}[1]
{
  \ifthenelse{\equal{#1}{0}}{non testé}{}
  \ifthenelse{\equal{#1}{1}}{\FiveStar\FiveStarOpen\FiveStarOpen\FiveStarOpen\FiveStarOpen}{}
  \ifthenelse{\equal{#1}{2}}{\FiveStar\FiveStar\FiveStarOpen\FiveStarOpen\FiveStarOpen}{}
  \ifthenelse{\equal{#1}{3}}{\FiveStar\FiveStar\FiveStar\FiveStarOpen\FiveStarOpen}{}
  \ifthenelse{\equal{#1}{4}}{\FiveStar\FiveStar\FiveStar\FiveStar\FiveStarOpen}{}
  \ifthenelse{\equal{#1}{5}}{\FiveStar\FiveStar\FiveStar\FiveStar\FiveStar}{}
}

\newcommand{\ingredient}[1][]{\ifthenelse{\equal{#1}{}}{\item }{\vspace{1ex}\hrule\vspace{1ex}\item[\textcurrency]\textbf{#1}}}

\newcommand{\etape}[1][]{\ifthenelse{\equal{#1}{}}{\item }{\item[\textbf{#1}]}}

% Environnement qui crée une nouvelle recette. Ingrédients, étape, cuisson etc doivent être contenus dans un environnement recette.
\newenvironment{recette}[4]{\newpage \nomRecette{#1}\statistiqueTopSep\note{#2}\ifthenelse{\equal{#4}{}}%
{\hfill\includegraphics[width=1em]{figures/logo_minuterie.pdf} \; \begin{itshape}\textbf{Préparation \& Cuisson :} #3\end{itshape}}%if none
{\hfill\includegraphics[width=1em]{figures/logo_minuterie.pdf} \; \begin{itshape}\textbf{Préparation :} #3 \hfill \includegraphics[width=1em]{figures/logo_minuterie.pdf} \;\textbf{Cuisson :} #4\end{itshape}}}{}

% Environnement pour la mise en page des ingrédients. Un argument optionnel de l'environnement permet de spécifier pour combien de personnes sont les doses. Les entrées de cet environnement doivent être ``\ingredient nom de l'ingrédient''. Si on veut séparer les ingrédients, pour dénoter deux sortes de choses, ingrédient pour une sauce et une pâte par exemple, il faut faire une ligne du style \ingredient[pour la sauce] afin de séparer.
\newenvironment{ingredients}[1][]{\ingredientsTopSep{\partStyle Ingrédients\ifthenelse{\equal{#1}{}}{}{ (#1)}}\begin{multicols}{2}\begin{itemize}\renewcommand{\labelitemi}{$\bullet$}}{\end{itemize}\end{multicols}}

% Environnement pour mettre en page la préparation du plat. Les entrées sont de la forme ``\etape texte de l'étape``. Si on veut séparer les étapes, pour dénoter un groupement d'étape par exemple, il faut faire une ligne du style \etape[pour la sauce] afin de séparer.
\newenvironment{preparation}{\preparationTopSep{\partStyle Préparation }\vspace{0.5em}\begin{enumerate}}{\end{enumerate}}

% À utiliser si on souhaite faire une mise en page un peu plus évoluée, avec deux listes par exemples.
\newenvironment{preparation*}{\preparationTopSep{\partStyle Préparation }\vspace{0.5em}\par}{}

% Environnement pour mettre en page la cuisson du plat
\newenvironment{cuisson}{\bigskip{\bfseries \large Cuisson }\par}{}
\makeatother

%TODO Faire une commande pour que dans l'index on voit la liste des recettes en fonction de la note qu'elles ont. 

\title{Recettes de Cuisine}
\author{Autiwa}
\begin{document}
\begin{titlepage}
\begin{center}
~
\vfill
% Upper part of the page
\begin{figure}[t]
\centering
\includegraphics[width=0.15\textwidth]{logo-autiwa.pdf}%chemin absolu de l'image pour que l'on puisse appeler ce fichier depuis plusieurs dossiers différents.
\end{figure}

% Title
\HRule \\[0.4cm]
{ \huge \bfseries \makeatletter\@title\makeatother}\\[0.4cm]

\HRule \\[0.75cm]
{\large \today}\\[0.75cm]
\makeatletter
\@author
\makeatother
\vfill
\vfill
~
% Bottom of the page


\end{center}
\end{titlepage}

\tableofcontents
\chapter{Recettes qui ne m'ont pas plu}
Il est possible d'essayer de les améliorer mais je n'ai généralement pas été fan du goût.

\begin{recette}{Gigot d'agneau rôti au lard}{3}{}{}
\begin{ingredients}
\ingredient gigot raccourci
\ingredient $150\unit{g}$ de fines tranches de poitrine fumée (prévoir le double si le gigot est relativement gros)
\ingredient 1 gousse d'ail
\ingredient 5 brins de romarins et de thym
\ingredient $50\unit{g}$ de beurre
\ingredient 1 cuillère à soupe de moutarde
\ingredient 1 cuillère à café de fond de veau déshydraté
\ingredient $10\unit{cl}$ de vin blanc sec
\ingredient sel,poivre
\end{ingredients}

\begin{preparation}
\etape Sortez le gigot du frigo 2h avant la cuisson. Allumez le fout th. 7 (210\degres C). Mélangez le beurre ramolli avec le thym effeuillé, le romarin et l'ail haché et étalez le sur le gigot.
\etape Posez le gigot dans un plat à rôtir et couvrez-le entièrement de tranches de poitrine fumée chevauchées. Glissez le plat dans le four. Laissez cuire $12\unit{min}$ par livre de viande (environ 1h).
\etape Retirez la viande cuite du plat et laissez-la reposer $20\unit{min}$ sous un papier d'alu. Dégraissez le jus, ajoutez la moutarde, le vin et le fond de veau dilué dans $15\unit{cl}$ d'eau.
\etape Mettez le plat sur le feu, faites bouillir $5\unit{min}$ en grattant pour décoller les sucs du fond. Versez en saucière. Tranchez le gigot et servez vite.
\end{preparation}

\begin{remarque}
Vous pouvez aussi ajouter dans le plat du gigot des gousses d'ails entières qui deviendront fondantes à l'issue de la cuisson.

Si le gigot est épais, il est possible qu'une heure ne soit pas suffisant pour le cuire. Surtout si le four n'a pas eu le temps de bien préchauffer.
\end{remarque}
\end{recette}

\begin{recette}{Soupe de poisson}{0}{}{}
\begin{ingredients}
\ingredient 500g de Lotte
\ingredient 500g de loup de mer
\ingredient 500g de rouget
\ingredient 500g de dorade
\ingredient 500g de coulis de tomate
\ingredient Un poireau
\ingredient 2 oignons
\ingredient 2 gousses d'ail
\ingredient sel, poivre, huile d'olive, sucre, jus de citron
\end{ingredients}

\begin{preparation}
\etape Hachez l'oignon et le poireau
\etape Dans une marmite haute, faites revenir l'oignon, l'ail et le poireau dans de l'huile d'olive. Laissez légèrement dorer le tout puis réservez-les.
\etape Faites dorer les poissons dans un peu d'huile, puis une fois colorés (pas besoin qu'ils soient cuits), ajoutez l'oignon et les poireaux pour laisser mijoter le tout à feux doux pendant 5 minutes
\etape Ajoutez le coulis de tomate et 2L d'eau et laissez bouillir 15 minutes. 
\etape Retirez les poissons et passez-les au presse purée (ou mixez les) et passez-les au tamis.
\etape Mettez la purée dans la marmite et complétez l'assaisonnement. 
\etape Rectifiez l'assaisonnement et portez à ébullition.
\end{preparation}

Servez très chaud accompagné de petits croûtons de pain aillé et de fromage rapé.
\end{recette}

\begin{recette}{Porc à la créole}{3}{}{}%\index{porc}\index{créole}
\begin{ingredients}
\ingredient $500 \unit{g}$ d'échine de porc
\ingredient $2$ tomates
\ingredient $2$ piments verts
\ingredient $1$ gousse d'ail
\ingredient $1$ oignon
\ingredient $1$ cuillère à soupe de thym émietté
\ingredient $1$ cuillère à soupe de curcuma (ou curry)
\ingredient $2$ cuillères à soupe d'huile
\end{ingredients}

\begin{preparation}
\etape Épluchez l'oignon et coupez-le en lamelles.
\etape Coupez la viande de porc en cubes.
\etape Épluchez les tomates : pour cela, plongez-les quelques minutes dans de l'eau bouillante et ôtez la peau à l'aide d'un couteau à lame fine.
\etape Coupez la chair des tomates en petits dés.
\etape Épluchez la gousse d'ail et les piments verts.
\etape Faites revenir l'oignon dans l'huile.
\etape Lorsqu'il a pris une belle couleur, ajoutez la viande de porc et laissez cuire pendant 3 minutes.
\etape Ajoutez le thym émietté, les dés de tomates, le curcuma (ou le curry), puis les piments et l'ail après les avoir écrasés.
\etape Laissez cuire le tout à feu doux pendant une dizaine de minutes.
\end{preparation}

\end{recette}

\begin{recette}{Paëlla (garniture déjà prête)}{0}{}{}
\begin{ingredients}
\ingredient garniture pour paëlla
\ingredient 2 cuillères à soupe d'huile
\ingredient $200$ g de riz
\ingredient $30$ cl d'eau
\ingredient épices pour paëlla
\end{ingredients}

\begin{preparation}
\etape Dans une grande poêle, faites chauffer à feu moyen l'huile.
\etape Versez le riz et laissez rissoler pendant 2 minutes environ en remuant de temps en temps.
\etape Versez l'eau dans la poêle et ajoutez les épices (safran et autres)
\etape Mélangez et portez à ébullition. Couvrez la poêle et laissez cuire 5 minutes jusqu'à absorption de l'eau.
\etape Rajoutez la garniture et répartissez son contenu sur le riz. Couvrez et faites cuire jusqu'à ce que le riz soit cuit et que le bouillon soit réduit. En ajoutant de l'eau si nécessaire.
\etape Servez dans la poêle de cuisson.
\end{preparation}

\end{recette}

\begin{recette}{Katlietkis}{2}{}{}
\begin{ingredients}
\ingredient 500g de viande hachée
\ingredient 250g de mie de pain
\ingredient 2 oignons
\ingredient 2 œufs
\ingredient sel aux herbes, poivre, aneth
\end{ingredients}

\begin{preparation}
\etape Faites ramolir la mie dans de l'eau pendant une petite heure. En gros, mettez la mie dans un saladier et mettez un peu d'eau.
\etape Égouttez la mie de pain (mettez la dans une passoire et appuyez sur la mie pour enlever le gros de l'eau).
\etape Mélangez la mie ainsi ramollie avec la viande, les oignons mixés (ou coupés très fin), les deux œufs. Poivrez, salez, rajoutez un peu d'aneth et mélangez bien.
\etape  Formez des boulettes (diamètre de 5--7 cm et épaisseur de quelques centimètres) puis passez les dans la chapelure

\etape[La préparation des galettes]

\etape Préparez de la chapelure dans une assiette.
\etape formez, d'une main (et avec une cuillère à soupe dans l'autre, une boule de garniture.
\etape posez là dans la chapelure sans l'écraser.
\etape avec la cuillère à soupe, saupoudrez de chapelure, puis retournez la.
\etape posez ensuite la boule dans la poele avec l'huile chaude, saisissez quelques minutes puis tournez afin que la galette prenne forme et ne se casse pas quand vous tournez avec la spatule.
\end{preparation}


\begin{cuisson}

Baissez le feu et laissez cuire 3/4 d'heure environ en les tournant de temps en temps. (Ne couvrez pas, afin que l'eau puisse s'évaporer.)

Si vous ne pouvez pas mettre toutes les galettes dans la poële en une seule fois, vous pouvez en faire cuire certaines au four une fois celles-ci dorées à la poële.

\begin{remarque}
Ça se garde quelques temps au frigo et ça se mange autant froid que chaud.
\end{remarque}
\end{cuisson}
\end{recette}

\end{document}